\chapter{模型方法与融合结构}

\section{输入表示}

本文使用两类输入表示：原始 I/Q 序列与 STFT 时频图。

\textbf{I/Q 序列。}基带采样信号记为 \(x_{\mathrm{IQ}}\in\mathbb{R}^{2\times128}\)，两行分别为同相分量 \(I[n]\) 和正交分量 \(Q[n]\)，\(n=0,\ldots,127\)。该表示直接保留时间采样结构，是大多数 AMC 文献的标准输入\cite{oshea2016cnn,sainath2015cldnn,xu2020mcldnn}。

\textbf{STFT 时频图。}以复数形式拼合 \(x_c[n]=I[n]+jQ[n]\)，对 \(x_c\) 施加窗函数 \(w[n]\) 后做短时傅里叶变换：
\begin{equation}
  X(\tau,k)=\sum_{n=0}^{N-1} x_c[n]\,w[n-\tau]\exp\!\left(-j2\pi kn/N\right),
  \label{eq:stft-definition}
\end{equation}
其中 \(\tau\) 为帧起点，\(k\) 为频率箱索引，\(N\) 为 FFT 点数。实验中选取 Hann 窗，窗长 32 点、跳步 16 点，FFT 大小 32，输出幅度谱形状为 \(\mathbb{R}^{1\times32\times7}\)（频率箱 \(\times\) 时间帧）。时频图以线性幅度馈入 2D-CNN，不取对数，以避免引入与信号幅度标定相关的额外超参。

\section{实验模型矩阵}

表~\ref{tab:method-family-matrix} 汇总本文实现的全部模型、输入视图、方法角色和参数量，对应第~\ref{sec:results-main}~节主实验的 9 个模型行及提案模型。

\begin{table}[htbp]
  \centering
  \caption{实验模型矩阵：输入视图、方法角色与参数量}
  \label{tab:method-family-matrix}
  \small
  \begin{tabularx}{\textwidth}{llXr}
    \toprule
    模型 & 输入视图 & 方法角色 & 参数量 \\
    \midrule
    \texttt{cnn1d}\cite{oshea2016cnn}         & I/Q      & 三层 Conv1d I/Q 基线，复现早期 CNN-AMC 思路 & 37,131 \\
    \texttt{resnet1d}\cite{he2016resnet}       & I/Q      & 残差连接深度卷积，提供更强 I/Q 对照 & 111,755 \\
    \texttt{tfcnn\_stft}                       & STFT     & 单独检验 STFT 视图的判别容量 & 24,123 \\
    \texttt{cldnn}\cite{sainath2015cldnn,west2017deep} & I/Q & 三层 Conv1d + 单层 LSTM 卷积-循环混合 & 241,675 \\
    \texttt{mcldnn}\cite{xu2020mcldnn}         & I/Q      & 多分支 2D/1D 卷积 + 2 层 LSTM 多通道时空结构 & 406,199 \\
    \texttt{lwamcnet}\cite{wang2021lwamcnet}   & I/Q      & 深度可分离残差卷积，参数量最低参考点 & 20,395 \\
    \texttt{iq\_param\_matched}                & I/Q      & 参数量对齐融合模型的纯 I/Q 控制 & 136,395 \\
    \texttt{fusion\_iq\_stft}                  & I/Q+STFT & 静态拼接双视图，检验时域-时频互补性 & 98,299 \\
    \texttt{gated\_fusion\_iq\_stft}           & I/Q+STFT & 标量门控自适应加权双视图 & 134,780 \\
    \midrule
    \texttt{fusion\_cldnn\_stft}（提案）       & I/Q+STFT & I/Q 分支升级为 CLDNN 风格，叠加信号域增强与标签平滑 & 286,000 \\
    \bottomrule
  \end{tabularx}
\end{table}

\section{I/Q 一维卷积基线}

\texttt{cnn1d} 以三层一维卷积堆叠从原始 I/Q 序列提取局部时间特征\cite{oshea2016cnn}。卷积核尺寸依次为 7、5、3，逐层扩大通道数（32-64-128）；第三层后接全局平均池化，消除对序列长度的依赖，最后以单层线性分类头输出 11 类 logits。表~\ref{tab:arch-cnn1d} 给出逐层结构。

\begin{table}[htbp]
  \centering
  \caption{\texttt{cnn1d} 逐层结构（输入 \([2,128]\)，batch 维省略）}
  \label{tab:arch-cnn1d}
  \small
  \begin{tabular}{llll}
    \toprule
    层 & 操作 & 输出形状 & 参数量 \\
    \midrule
    Conv1d-1  & kernel 7, padding 3, 32ch, BN, ReLU  & \([32,128]\) & 480 \\
    MaxPool1d & kernel 2                              & \([32, 64]\) & --- \\
    Conv1d-2  & kernel 5, padding 2, 64ch, BN, ReLU  & \([64, 64]\) & 10,304 \\
    MaxPool1d & kernel 2                              & \([64, 32]\) & --- \\
    Conv1d-3  & kernel 3, padding 1, 128ch, BN, ReLU & \([128,32]\) & 24,832 \\
    AvgPool   & AdaptiveAvgPool1d(1)                  & \([128]\)    & --- \\
    Linear    & \(128 \to 11\)                        & \([11]\)     & 1,419 \\
    \bottomrule
  \end{tabular}
\end{table}

\section{I/Q 残差卷积：\texttt{resnet1d}}

\texttt{resnet1d} 在一维卷积堆叠中引入残差连接\cite{he2016resnet}，允许堆叠更深特征层。主干由一个 stem 卷积层和三个残差块组成，Block2、Block3 以 stride=2 完成下采样；各残差块在通道数或步长不匹配时用 \(1\times1\) 卷积对齐快捷路径，否则退化为恒等映射。全局平均池化后接线性分类头，共 111,755 个参数（表~\ref{tab:arch-resnet1d}）。

\begin{table}[htbp]
  \centering
  \caption{\texttt{resnet1d} 逐层结构（输入 \([2,128]\)，batch 维省略）}
  \label{tab:arch-resnet1d}
  \small
  \begin{tabular}{llll}
    \toprule
    层 & 操作 & 输出形状 & 说明 \\
    \midrule
    Stem      & Conv1d(2\(\to\)32, k=7, p=3), BN, ReLU  & \([32,128]\) & --- \\
    Block1    & ResBlock(32\(\to\)32, stride=1)           & \([32,128]\) & 恒等快捷 \\
    Block2    & ResBlock(32\(\to\)64, stride=2)           & \([64, 64]\) & 1{\texttimes}1 快捷 \\
    Block3    & ResBlock(64\(\to\)128, stride=2)          & \([128,32]\) & 1{\texttimes}1 快捷 \\
    AvgPool   & AdaptiveAvgPool1d(1), flatten             & \([128]\)    & --- \\
    Linear    & \(128 \to 11\)                            & \([11]\)     & --- \\
    \bottomrule
  \end{tabular}
\end{table}
\noindent 各 ResBlock 由两层 Conv1d(k=3, p=1)+BN 构成，中间接 ReLU，快捷路径在块输出处求和后再经 ReLU。

\section{卷积-循环混合模型：\texttt{cldnn} 与 \texttt{mcldnn}}

卷积-循环堆叠（CLDNN）最早在语音识别中系统描述\cite{sainath2015cldnn}：先以卷积层压缩局部变化并提取空间特征，再以 LSTM\cite{hochreiter1997lstm} 建模帧间时序依赖，最后以全连接层分类。该思路被引入 AMC 比较研究\cite{west2017deep}，与纯卷积基线形成结构对照。

\subsection{\texttt{cldnn}}

本文 \texttt{cldnn} 为单流单层 LSTM 结构（表~\ref{tab:arch-cldnn}）。I/Q 序列经三层 Conv1d 提取局部特征，前两层后接 MaxPool1d 各做 2 倍下采样；第三层无池化，将特征序列转置为 LSTM 输入 \([B,32,128]\)；单层 LSTM 取末时刻隐状态 \(h_T\in\mathbb{R}^{128}\)，送入两层 MLP 分类器（Dropout 0.5）。

\begin{table}[htbp]
  \centering
  \caption{\texttt{cldnn} 逐层结构（输入 \([2,128]\)，batch 维省略）}
  \label{tab:arch-cldnn}
  \small
  \begin{tabular}{llll}
    \toprule
    层 & 操作 & 输出形状 & 说明 \\
    \midrule
    Conv1d-1  & kernel 7, padding 3, 64ch, BN, ReLU  & \([64,128]\) & --- \\
    MaxPool1d & kernel 2                               & \([64, 64]\) & --- \\
    Conv1d-2  & kernel 5, padding 2, 128ch, BN, ReLU & \([128,64]\) & --- \\
    MaxPool1d & kernel 2                               & \([128,32]\) & --- \\
    Conv1d-3  & kernel 3, padding 1, 128ch, BN, ReLU & \([128,32]\) & --- \\
    转置      & transpose(1,2)                         & \([32,128]\) & LSTM 输入 \\
    LSTM      & hidden=128, 1 层, batch\_first，取 \(h_T\) & \([128]\) & --- \\
    Linear-1  & \(128\to128\), ReLU, Dropout(0.5)    & \([128]\)    & --- \\
    Linear-2  & \(128\to11\)                          & \([11]\)     & --- \\
    \bottomrule
  \end{tabular}
\end{table}

\subsection{\texttt{mcldnn}}

MCLDNN 引入多分支 I/Q 空间特征提取\cite{xu2020mcldnn}：以 \((2,8)\) 核对联合 I/Q 矩阵做 2D 卷积提取联合特征，同时对 I、Q 各做因果 1D 卷积后经 \((1,8)\) 核做空间融合；两路拼接后经 \((2,5)\) 融合卷积压缩至时间序列，馈入 2 层 LSTM（hidden=128）；分类器为三层线性-SELU-Dropout 结构（Dropout 0.5）。表~\ref{tab:arch-mcldnn} 给出逐层结构。

\begin{table}[htbp]
  \centering
  \caption{\texttt{mcldnn} 逐层结构（输入 \([2,128]\)，展开为 \([1,2,128]\) 馈入 2D 分支，batch 维省略）}
  \label{tab:arch-mcldnn}
  \small
  \begin{tabular}{llll}
    \toprule
    分支 / 层 & 操作 & 输出形状 & 说明 \\
    \midrule
    \multicolumn{4}{l}{\textit{联合空间分支}} \\
    \quad Conv2d-joint & SamePad(1\(\to\)50, k=(2,8)), ReLU  & \([50,1,128]\) & 联合 I/Q 2D \\
    \multicolumn{4}{l}{\textit{独立分支（I、Q 各一路）}} \\
    \quad CausalConv-I & CausalConv(1\(\to\)50, k=8), ReLU   & \([50,128]\)   & I 单通道 \\
    \quad CausalConv-Q & CausalConv(1\(\to\)50, k=8), ReLU   & \([50,128]\)   & Q 单通道 \\
    \quad Stack I/Q    & stack(I,Q, dim=2)                    & \([50,2,128]\) & --- \\
    \quad Conv2d-indep & SamePad(50\(\to\)50, k=(1,8)), ReLU & \([50,2,128]\) & --- \\
    \multicolumn{4}{l}{\textit{融合阶段}} \\
    \quad Concat        & 联合+独立，dim=1                    & \([100,2,128]\) & --- \\
    \quad Conv2d-fuse   & (100\(\to\)100, k=(2,5)), ReLU      & \([100,1,124]\) & --- \\
    \quad Squeeze+转置  & squeeze(2), transpose(1,2)          & \([124,100]\)  & LSTM 输入 \\
    \quad LSTM          & hidden=128, 2 层, 取 \(h_T\)       & \([128]\)      & --- \\
    \quad Linear-1      & \(128\to128\), SELU, Dropout(0.5)  & \([128]\)      & --- \\
    \quad Linear-2      & \(128\to128\), SELU, Dropout(0.5)  & \([128]\)      & --- \\
    \quad Linear-3      & \(128\to11\)                        & \([11]\)       & --- \\
    \bottomrule
  \end{tabular}
\end{table}
\noindent\textbf{MCLDNN 训练异常说明。}本实验中 \texttt{mcldnn} 三个训练种子表现出双峰分布：种子 42 接近正常，种子 2025 与 3407 退化至接近随机水平，导致整体准确率均值与标准差均不稳定（见第~\ref{sec:mcldnn-anomaly}~节）。该异常作为负面证据完整保留，不从结果表中剔除。

\section{轻量深度可分离卷积：\texttt{lwamcnet}}

\texttt{lwamcnet} 对应文献\cite{wang2021lwamcnet} 的 2016.10A 缩放方案：以 \((2,5)\) 核提取初始 I/Q 空间特征，经三级深度可分离残差（DSC-Residual）栈精炼，每级含两个 DSC 残差单元和 \((1,2)\) MaxPool；最后以分组深度卷积（GDWConv，k=(1,16)）重建特征，接单层线性分类器。深度可分离卷积将标准 2D 卷积分解为逐通道卷积（depthwise）和逐点卷积（pointwise），在保持感受野的同时大幅降低参数量，本文 20,395 个参数是所有模型中最少的（表~\ref{tab:arch-lwamcnet}）。

\begin{table}[htbp]
  \centering
  \caption{\texttt{lwamcnet} 逐层结构（输入 \([2,128]\)，以 \([1,2,128]\) 馈入首层，batch 维省略）}
  \label{tab:arch-lwamcnet}
  \small
  \begin{tabular}{llll}
    \toprule
    层 & 操作 & 输出形状 & 说明 \\
    \midrule
    Conv2d-0    & (1\(\to\)64, k=(2,5), p=(0,2)), ReLU & \([64,1,128]\) & 初始特征 \\
    DSC-Stack-1 & 64\(\to\)32: ChannelFusion+2{\texttimes}DSCUnit+MaxPool(1,2) & \([32,1,64]\) & --- \\
    DSC-Stack-2 & 32\(\to\)32: ChannelFusion+2{\texttimes}DSCUnit+MaxPool(1,2) & \([32,1,32]\) & --- \\
    DSC-Stack-3 & 32\(\to\)32: ChannelFusion+2{\texttimes}DSCUnit+MaxPool(1,2) & \([32,1,16]\) & --- \\
    GDWConv     & groups=32, k=(1,16), ReLU, flatten   & \([32]\) & --- \\
    Linear      & \(32\to11\)                           & \([11]\) & --- \\
    \bottomrule
  \end{tabular}
\end{table}
\noindent 各 DSCUnit 含两个深度可分离 Conv2d（k=(1,5), p=(0,2)）串联作残差路径，等效于对时间维做深度序列卷积。

\section{STFT 单视图模型：\texttt{tfcnn\_stft}}

\texttt{tfcnn\_stft} 仅以 STFT 幅度谱 \(x_{\mathrm{STFT}}\in\mathbb{R}^{1\times32\times7}\) 为输入，通过三层 2D-CNN 提取时频局部模式，经全局平均池化后以线性分类头输出预测。该模型单独度量 STFT 视图的判别容量，为解释双视图融合效果提供孤立视图参照（表~\ref{tab:arch-tfcnn}）。

\begin{table}[htbp]
  \centering
  \caption{\texttt{tfcnn\_stft} 逐层结构（输入 \([1,32,7]\)，batch 维省略）}
  \label{tab:arch-tfcnn}
  \small
  \begin{tabular}{llll}
    \toprule
    层 & 操作 & 输出形状 & 说明 \\
    \midrule
    Conv2d-1  & kernel 3, padding 1, 16ch, BN, ReLU & \([16,32,7]\) & --- \\
    MaxPool2d & kernel 2                              & \([16,16,3]\) & --- \\
    Conv2d-2  & kernel 3, padding 1, 32ch, BN, ReLU & \([32,16,3]\) & --- \\
    MaxPool2d & kernel 2                              & \([32, 8,1]\) & --- \\
    Conv2d-3  & kernel 3, padding 1, 64ch, BN, ReLU & \([64, 8,1]\) & --- \\
    AvgPool   & AdaptiveAvgPool2d(1), flatten         & \([64]\)      & --- \\
    Linear    & \(64 \to 11\)                         & \([11]\)      & --- \\
    \bottomrule
  \end{tabular}
\end{table}

\section{参数匹配 I/Q 控制：\texttt{iq\_param\_matched}}

双视图融合相比纯 I/Q 模型不仅引入额外视图，还增加 STFT 分支参数。为区分"额外参数"与"新增视图信息"的贡献，本文引入 \texttt{iq\_param\_matched}：仅以 I/Q 为输入，通过四层 Conv1d（通道数 64-96-128-128）使参数量（136,395）与 \texttt{gated\_fusion\_iq\_stft}（134,780）基本对齐。当 \texttt{fusion\_iq\_stft} 在低信噪比区间相对该控制有小幅正差、整体区间为负差时，可将差异解释为视图信息而非参数规模的效果（表~\ref{tab:arch-iqpm}）。

\begin{table}[htbp]
  \centering
  \caption{\texttt{iq\_param\_matched} 逐层结构（输入 \([2,128]\)，batch 维省略）}
  \label{tab:arch-iqpm}
  \small
  \begin{tabular}{llll}
    \toprule
    层 & 操作 & 输出形状 & 说明 \\
    \midrule
    Conv1d-1  & kernel 7, padding 3, 64ch, BN, ReLU  & \([64,128]\) & --- \\
    MaxPool1d & kernel 2                               & \([64, 64]\) & --- \\
    Conv1d-2  & kernel 5, padding 2, 96ch, BN, ReLU  & \([96, 64]\) & --- \\
    MaxPool1d & kernel 2                               & \([96, 32]\) & --- \\
    Conv1d-3  & kernel 3, padding 1, 128ch, BN, ReLU & \([128,32]\) & --- \\
    Conv1d-4  & kernel 3, padding 1, 128ch, BN, ReLU & \([128,32]\) & --- \\
    AvgPool   & AdaptiveAvgPool1d(1), flatten          & \([128]\)    & --- \\
    Linear-1  & \(128\to128\), ReLU, Dropout(0.1)    & \([128]\)    & --- \\
    Linear-2  & \(128\to11\)                          & \([11]\)     & --- \\
    \bottomrule
  \end{tabular}
\end{table}

\section{静态 I/Q--STFT 融合：\texttt{fusion\_iq\_stft}}\label{sec:method-static-fusion}

\texttt{fusion\_iq\_stft} 包含一个 I/Q 分支（\texttt{IQBranch1D}，三层 Conv1d，结构与 \texttt{cnn1d} 相同，输出 128 维嵌入）和一个 STFT 分支（\texttt{TimeFrequencyBranch2D}，三层 Conv2d，结构与 \texttt{tfcnn\_stft} 相同，输出 64 维嵌入），两分支独立提取特征后拼接经 MLP 分类：
\begin{equation}
  z = \bigl[h_{\mathrm{IQ}}\;\|\;h_{\mathrm{STFT}}\bigr]\in\mathbb{R}^{192},\qquad
  \hat{y} = \mathrm{MLP}_2(z),
  \label{eq:static-fusion}
\end{equation}
MLP 分类头结构为 \(192\to192\)（ReLU, Dropout 0.1）\(\to11\)。该设计检验 STFT 时频特征与 I/Q 时域特征能否形成互补；第~\ref{sec:results-extended}~节的乘加量分析（I/Q 分支 2.0 M MACs vs.\ \texttt{resnet1d} 4.9 M）表明，I/Q 分支表征容量是限制融合性能的主要瓶颈。

\section{标量门控融合：\texttt{gated\_fusion\_iq\_stft}}

\texttt{gated\_fusion\_iq\_stft} 在双分支基础上增加可学习标量门控。两分支嵌入经投影层映射至公共维度 \(d=128\)（\(W_{\mathrm{IQ}}\in\mathbb{R}^{128\times128}\)，\(W_{\mathrm{STFT}}\in\mathbb{R}^{128\times64}\)），拼接绝对差与逐元素乘积构成门控输入，经小型 MLP 输出标量权重 \(g\in(0,1)\)：
\begin{equation}
  r=\bigl[z_{\mathrm{IQ}}\;\|\;z_{\mathrm{STFT}}\;\|\;|z_{\mathrm{IQ}}-z_{\mathrm{STFT}}|\;\|\;z_{\mathrm{IQ}}\odot z_{\mathrm{STFT}}\bigr]\in\mathbb{R}^{512},
  \label{eq:gated-input}
\end{equation}
\begin{equation}
  g=\sigma\!\bigl(W_2\,\mathrm{ReLU}(W_1 r)\bigr),\qquad
  h=(1-g)\,z_{\mathrm{IQ}}+g\,z_{\mathrm{STFT}}\in\mathbb{R}^{128},
  \label{eq:gated-fusion}
\end{equation}
其中 \(W_1\in\mathbb{R}^{64\times512}\)，\(W_2\in\mathbb{R}^{1\times64}\)，\(\mathrm{ReLU}\) 后施加 Dropout(0.1)，\(\sigma\) 为 sigmoid。融合向量 \(h\) 经 \(128\to128\to11\) 分类头输出 logits。门控输入加入绝对差与逐元素乘积，使 MLP 能感知两视图间的差异幅度与方向，从而做出样本级视图加权；该门控不接受显式信噪比输入，是自适应加权的可检验候选。

\begin{figure}[htbp]
  \centering
  \resizebox{\textwidth}{!}{%
\begin{tikzpicture}[
  block/.style={rounded corners, draw=#1!75!black, fill=#1!10, very thick,
    align=center, minimum width=2.45cm, minimum height=0.9cm, font=\small},
  smallblock/.style={rounded corners, draw=black!55, fill=black!3, thick,
    align=center, minimum width=2.4cm, minimum height=0.75cm, font=\scriptsize},
  arrow/.style={-{Latex[length=2.1mm]}, thick, draw=black!65}
]
\node[block=blue] (iqin) at (0,1.35) {I/Q input\\$x_{\mathrm{IQ}}$};
\node[smallblock] (iqbranch) at (2.9,1.35) {1D Conv branch\\$h_{\mathrm{IQ}}$};
\node[block=teal] (stftin) at (0,-1.35) {STFT input\\$x_{\mathrm{STFT}}$};
\node[smallblock] (stftbranch) at (2.9,-1.35) {2D Conv branch\\$h_{\mathrm{STFT}}$};
\node[smallblock] (proj1) at (5.6,1.35) {Projection\\$z_{\mathrm{IQ}}$};
\node[smallblock] (proj2) at (5.6,-1.35) {Projection\\$z_{\mathrm{STFT}}$};
\node[block=orange] (concat) at (8.7,1.35) {Static fusion\\concat + classifier};
\node[block=purple] (gate) at (8.7,-1.35) {Scalar gate\\$(1-g)z_{\mathrm{IQ}}+gz_{\mathrm{STFT}}$};
\node[smallblock] (out1) at (11.8,1.35) {Class logits};
\node[smallblock] (out2) at (11.8,-1.35) {Class logits};
\draw[arrow] (iqin) -- (iqbranch);
\draw[arrow] (stftin) -- (stftbranch);
\draw[arrow] (iqbranch) -- (proj1);
\draw[arrow] (stftbranch) -- (proj2);
\draw[arrow] (proj1) -- (concat);
\draw[arrow] (proj2) -- (concat);
\draw[arrow] (concat) -- (out1);
\draw[arrow] (proj1) -- (gate);
\draw[arrow] (proj2) -- (gate);
\draw[arrow] (gate) -- (out2);
\node[draw=black!45, rounded corners, fill=gray!7, align=center, font=\scriptsize,
  minimum width=11.7cm] at (5.95,-2.75)
  {The gate is a tested candidate, not a supported robustness claim; Stage 5A/5B results decide its interpretation.};
\end{tikzpicture}%
}

  \caption{静态 I/Q--STFT 融合与标量门控融合结构对比。静态融合直接拼接两分支嵌入（\(\mathbb{R}^{128+64}\)）；门控融合先投影至公共维度（\(\mathbb{R}^{128}\)），再以样本级标量权重组合。}
  \label{fig:fusion-architecture}
\end{figure}

\section{提案模型与训练侧干预}\label{sec:method-proposed}

\subsection{架构动机}

第~\ref{sec:results-extended}~节乘加量分析显示，\texttt{fusion\_iq\_stft} 的 I/Q 分支 MACs 仅 2.0 M，而 \texttt{resnet1d} 为 4.9 M；换言之，``融合输给 CLDNN'' 更可能源于 I/Q 表征容量不足，而非多视图结构本身的局限。基于此，提案模型 \texttt{fusion\_cldnn\_stft} 将 I/Q 分支替换为 \texttt{CLDNNIQBranch}——与 \texttt{cldnn} 共享卷积特征栈（Conv1d 64-128-128ch）加单层 LSTM\cite{sainath2015cldnn,hochreiter1997lstm}，取末时刻隐状态 \(h_T\) 为分支嵌入；STFT 分支沿用 \texttt{TimeFrequencyBranch2D}（与 \texttt{tfcnn\_stft} 结构相同\cite{ioffe2015batchnorm}），分类头为小型 MLP\cite{srivastava2014dropout}：
\begin{equation}
  h_{\mathrm{IQ}}=\mathrm{LSTM}\bigl(f_{\mathrm{1D}}(x_{\mathrm{IQ}})\bigr)_T,\quad
  h_{\mathrm{STFT}}=f_{\mathrm{2D}}(x_{\mathrm{STFT}}),\quad
  \hat{y}=\mathrm{MLP}_2\!\bigl([h_{\mathrm{IQ}}\,\|\,h_{\mathrm{STFT}}]\bigr).
  \label{eq:fusion-cldnn-stft}
\end{equation}
模型总参数量 286 k（介于 \texttt{cldnn} 的 241.7 k 与 \texttt{mcldnn} 的 406.2 k 之间），逐层结构见表~\ref{tab:fusion-cldnn-stft-arch}。

\begin{table}[htbp]
  \centering
  \caption{提案模型 \texttt{fusion\_cldnn\_stft} 逐层结构（batch=1，形状按各分支独立标记）}
  \label{tab:fusion-cldnn-stft-arch}
  \small
  \begin{tabular}{llll}
    \toprule
    分支 / 层 & 操作 & 输出形状 & 说明 \\
    \midrule
    \multicolumn{4}{l}{\textit{I/Q 分支 (\texttt{CLDNNIQBranch})}} \\
    \quad Conv1d-1  & kernel 7, padding 3, 64ch, BN, ReLU  & \([64,128]\) & --- \\
    \quad MaxPool1d & kernel 2                               & \([64, 64]\) & --- \\
    \quad Conv1d-2  & kernel 5, padding 2, 128ch, BN, ReLU & \([128,64]\) & --- \\
    \quad MaxPool1d & kernel 2                               & \([128,32]\) & --- \\
    \quad Conv1d-3  & kernel 3, padding 1, 128ch, BN, ReLU & \([128,32]\) & --- \\
    \quad LSTM      & hidden=128, 1 层, 取 \(h_T\)         & \([128]\)    & batch\_first \\
    \midrule
    \multicolumn{4}{l}{\textit{STFT 分支 (\texttt{TimeFrequencyBranch2D})}} \\
    \quad Conv2d-1  & kernel 3, padding 1, 16ch, BN, ReLU  & \([16,32,7]\) & --- \\
    \quad MaxPool2d & kernel 2                               & \([16,16,3]\) & --- \\
    \quad Conv2d-2  & kernel 3, padding 1, 32ch, BN, ReLU  & \([32,16,3]\) & --- \\
    \quad MaxPool2d & kernel 2                               & \([32, 8,1]\) & --- \\
    \quad Conv2d-3  & kernel 3, padding 1, 64ch, BN, ReLU  & \([64, 8,1]\) & --- \\
    \quad AvgPool   & AdaptiveAvgPool2d(1), flatten          & \([64]\)      & --- \\
    \midrule
    \multicolumn{4}{l}{\textit{融合分类头}} \\
    \quad Concat    & \([128]\,\|\,[64]\)                   & \([192]\)    & --- \\
    \quad Linear-1  & \(192\to192\), ReLU, Dropout(0.1)    & \([192]\)    & --- \\
    \quad Linear-2  & \(192\to11\)                          & \([11]\)     & --- \\
    \bottomrule
  \end{tabular}
\end{table}

\subsection{信号域增强}

信号域增强仅在训练集启用，以保留 I/Q 物理可解释性。本文采用两种变换：

\textbf{随机相位旋转。}以概率 0.5 对每个训练样本施加角度 \(\theta\sim\mathcal{U}[0,2\pi)\) 的相位旋转：
\begin{equation}
  \begin{bmatrix}I'[n]\\Q'[n]\end{bmatrix}
  =\begin{bmatrix}\cos\theta & -\sin\theta\\\sin\theta & \cos\theta\end{bmatrix}
  \begin{bmatrix}I[n]\\Q[n]\end{bmatrix},\quad n=0,\ldots,127.
  \label{eq:phase-rotation}
\end{equation}
相位旋转保持各调制信号的星座图结构，不改变有效 SNR，分类标签不变。

\textbf{随机循环时移。}以概率 0.5 对样本施加随机整数时移 \(k\sim\mathcal{U}\{-32,\ldots,+32\}\)：
\begin{equation}
  I'[n]=I[(n-k)\bmod 128],\quad Q'[n]=Q[(n-k)\bmod 128].
  \label{eq:cyclic-shift}
\end{equation}
循环时移扩充序列对齐多样性，不影响频域特征分布，有效 SNR 与分类标签不变。

幅度缩放与加性高斯噪声因改变有效 SNR 而不予采用；图像式 mixup 因对 I/Q 线性插值缺乏物理可解释性\cite{huang2022mixingsignals,shorten2019dataaug}也不采用。图~\ref{fig:augmentation-demo} 给出四种情形的波形可视化对比。

\subsection{标签平滑}

标签平滑 \(\epsilon=0.1\)\cite{szegedy2016inception} 将 one-hot 目标 \(\mathbf{1}_{y}\) 替换为软目标分布：
\begin{equation}
  q_k=\begin{cases}1-\varepsilon+\varepsilon/K & k=y,\\
  \varepsilon/K & k\neq y,\end{cases}
  \quad K=11.
  \label{eq:label-smoothing}
\end{equation}
该操作抑制过度自信，是分类任务的常见正则化手段，在本实验中由 \texttt{train.loss} 字段配置。

\subsection{三项干预的不平等贡献}\label{sec:method-ablation-preview}

三项干预对整体准确率的贡献并不对等：消融实验（第~\ref{sec:proposed-ablation}~节）显示，架构升级本身对整体准确率的增益约 0 pp（相对 \texttt{cldnn} 仅 \(-0.13\) pp，在种子噪声范围内）；信号域增强是主要贡献项，净增益约 \(+1.55\) pp（arch+aug vs. arch\_only，McNemar \(p\approx6\times10^{-78}\)）；标签平滑则略有负效（\(-0.19\) pp，95\% CI $[-0.0033,-0.0006]$，\(p=0.0155\)）。因此，提案模型的实质贡献应理解为"验证信号域增强在卷积-循环融合架构上的有效性"，架构升级与标签平滑均非主要来源。

\begin{figure}[htbp]
  \centering
  \includegraphics[width=0.96\textwidth]{fig23_augmentation_visualization.pdf}
  \caption{信号域增强示例（QPSK @ +6 dB 单样本）。(a) 原始；(b) 相位旋转 $\theta=\pi/3$；(c) 循环时移 $k=+8$；(d) 两者叠加。I/Q 通道线型见图例。}
  \label{fig:augmentation-demo}
\end{figure}

\section{实现细节与比较纪律}\label{sec:method-impl}

所有模型使用 PyTorch~2.9.1+cu128\cite{paszke2019pytorch} 与 NVIDIA driver 570.\(\ast\) 实现，数据集固定划分（训练/验证/测试集一次性确定，全程不变）。主实验 9 个模型行在 RTX 4070 Ti (Ada, sm\_89) 上训练；提案模型及消融变体在 RTX 3090 (Ampere, sm\_86) 上训练。两批实验的硬件差异作为协议混淆因素显式登记，不参与跨批直接数值比较。

共用训练设置：优化器 AdamW\cite{loshchilov2019adamw,kingma2015adam}，学习率 \(\eta=10^{-3}\)（AMC 卷积网络的常用设定\cite{xu2020mcldnn,zhang2021petcgdnn}），权重衰减 \(10^{-4}\)\cite{krogh1992weightdecay}，批量大小 256。主实验矩阵最大训练轮数 30，早停耐心值 8（依据验证集准确率）；提案模型及消融变体最大轮数 50，早停耐心值 15，额外采用 5 轮线性 warmup 加 cosine annealing 学习率调度\cite{loshchilov2017sgdr}。损失函数默认为标准交叉熵；提案模型加 \(\varepsilon=0.1\) 标签平滑（式~\ref{eq:label-smoothing}）。所有模型以种子集合 \{42, 2025, 3407\} 各训练一次，报告均值 \(\pm\) 标准差。

统一超参作用于主实验矩阵全部 9 个模型，不做模型级事后调参，以保持比较纪律。提案模型及消融变体独立评估，结论限于第~\ref{sec:proposed-ablation}~节范围内，不与主实验矩阵做直接数值比较。本章方法描述限定于已登记模型与训练设置，不引入未在第~4~章评估的诊断性候选。
